\documentclass[11pt,a4paper]{article}
\usepackage[margin=25mm,headheight=16pt]{geometry}
\usepackage{fontspec,kotex,xcolor,hyperref,fancyhdr}
\setmainfont{DejaVu Serif}
\setmainhangulfont{Noto Sans CJK KR}
\definecolor{navy}{HTML}{173E59}
\pagestyle{fancy}\fancyhf{}
\fancyhead[L]{\small GuardSynth CoC}
\fancyhead[R]{\small Title and abstract draft · English}
\fancyfoot[L]{\small v0.1 · 2026-09-30}
\fancyfoot[R]{\small\thepage}
\renewcommand{\headrulewidth}{0.3pt}
\setlength{\parindent}{0pt}
\setlength{\emergencystretch}{3em}
\linespread{1.15}
\hypersetup{pdftitle={GuardSynth: Formal Specification and Verification of Behavioral Constraints for CoC-Based VLM Learning},pdfsubject={Title and abstract for author review},pdfauthor={GuardSynth project}}
\begin{document}
\vspace*{8mm}
{\Large\bfseries\color{navy} GuardSynth: Formal Specification and Verification of Behavioral Constraints for CoC-Based VLM Learning\par}
\vspace{8mm}
{\large\bfseries Abstract\par}
\vspace{3mm}
Chain of Cognition (CoC) connects scene understanding with action reasoning for vision-language model (VLM) learning. However, learning from CoC alone can yield poor action selection when constraint applicability and release conditions remain unspecified. We present GuardSynth, a framework that augments CoC with explicit behavioral constraints. We define their scope and constituent elements and specify applicability, prohibited actions, persistence and release conditions, timing, and evidence using Evidence-Bound Lifecycle Contracts (EBLC). We then verify consistency and designated properties within a restricted logical model, generate Controlled Natural Language (CNL) corresponding to specification fields, and add it to CoC training inputs. We compare Qwen3-VL-2B-Instruct under matched LoRA training budgets across three seeds on a controlled synthetic visual-temporal task. Mean action-selection accuracy is 49.86\% with CoC alone, 96.81\% with our existing natural-language constraints, and 98.06\% with EBLC-derived CNL; corresponding constraint-violation rates are 19.17\%, 1.39\%, and 0.83\%. These results demonstrate a systematic path from behavioral constraint specification and verification to CoC augmentation, alongside improved action selection with explicit constraints. The findings concern a controlled task with constraints supplied during both training and evaluation; they establish neither a causal performance benefit of verification itself nor real-road safety.
\end{document}
